% GENERATED FILE. Edit canonical lessons, then run npm run generate:books.
% canonical-source-hash: fnv1a64:9a752a59b4cb412f
% canonical-lessons: SA-C33-river, SA-C33-bridge, SA-C33-mountain, SA-C33-town, SA-C33-sea

\chapter{The River and the Road}
\label{ch:sa-land}
\hlchaptermodality{33}

\begin{chapteropening}
\textbf{By the end of this chapter:} \emph{I can name a river, a bridge, a mountain, a town and the sea in Sanskrit, and say which of them is a cousin of the Greek word behind \emph{police}.}

Complete SA-C33-sea and demonstrate the chapter promise.
\end{chapteropening}

\section[nadī]{\sk{नदी} (nadī) — a river}

\label{lesson:SA-C33-river}



\begin{quote}
Before the new word: what did \emph{dāru} mean?
\end{quote}

\subsection*{You'll want to know: \sk{नदी}}
\textbf{\sk{नदी}} (\emph{nadī}) — \textquotedblleft{}a river\textquotedblright{}.

\textbf{\sk{नदी}} says what a river does rather than what it is: the root behind it means \textquotedblleft{}to roar, to sound\textquotedblright{}. A river is the noisy one.

It ends in long \emph{-ī}, like \textbf{\sk{नारी}} and \textbf{\sk{भगिनी}}, and it has outlasted almost everything around it — Hindi, Marathi, Bengali, Gujarati and Nepali all still say \emph{nadī}, and every river name ending in \emph{-nadi} on an Indian map carries it whole.

\begin{grammarlens}[title={What you've built}]
A river, and a word that survived every one of its neighbours.
\end{grammarlens}

\subsection*{Your turn}
Say these aloud:

\begin{itemize}
\raggedright
  \item \emph{nadī}
  \item it once more, slowly
  \item \emph{dāru}, then \emph{nadī}
\end{itemize}

\begin{tcolorbox}[breakable,colback=teal!4,colframe=teal!35!black,title={Before you move on}]
What does \sk{नदी} mean? (\textquotedblleft{}a river\textquotedblright{}.) Say it once more.
\end{tcolorbox}
\section[setuḥ]{\sk{सेतुः} (setuḥ) — a bridge, a causeway}

\label{lesson:SA-C33-bridge}



\begin{quote}
Before the new word: what did \emph{nadī} mean?
\end{quote}

\subsection*{You'll want to know: \sk{सेतुः}}
\textbf{\sk{सेतुः}} (\emph{setuḥ}) — \textquotedblleft{}a bridge, a causeway\textquotedblright{}.

\textbf{\sk{सेतुः}} is literally \textquotedblleft{}the binding\textquotedblright{}, from a root meaning \textquotedblleft{}to bind, to tie\textquotedblright{}. A bridge is what ties two banks to each other, and Sanskrit names it for the tying rather than for the crossing.

The causeway between India and Lanka is the \emph{Rāma-setu}, and that is this word, unchanged, on a modern map. Ends in \emph{-uḥ}, like \textbf{\sk{तरुः}}, \textbf{\sk{वायुः}} and \textbf{\sk{गुरुः}}.

\begin{grammarlens}[title={What you've built}]
A bridge, and a word that means the tying rather than the crossing.
\end{grammarlens}

\subsection*{Your turn}
Say these aloud:

\begin{itemize}
\raggedright
  \item \emph{setuḥ}
  \item it once more, slowly
  \item \emph{nadī}, then \emph{setuḥ}
\end{itemize}

\begin{tcolorbox}[breakable,colback=teal!4,colframe=teal!35!black,title={Before you move on}]
What does \sk{सेतुः} mean? (\textquotedblleft{}a bridge, a causeway\textquotedblright{}.) Say it once more.
\end{tcolorbox}
\section[parvataḥ]{\sk{पर्वतः} (parvataḥ) — a mountain}

\label{lesson:SA-C33-mountain}



\begin{quote}
Before the new word: what did \emph{setuḥ} mean?
\end{quote}

\subsection*{You'll want to know: \sk{पर्वतः}}
\textbf{\sk{पर्वतः}} (\emph{parvataḥ}) — \textquotedblleft{}a mountain\textquotedblright{}.

\textbf{\sk{पर्वतः}} is built on \emph{parvan}, \textquotedblleft{}a knot, a joint\textquotedblright{} — the same word for the joints of a bamboo cane and for the sections of a long poem. A mountain is the knotted, jointed land, and a range of them is a knotted spine.

Hindi, Marathi and Nepali keep \emph{parvat}. You met one range of them last chapter, in the snow.

\begin{grammarlens}[title={What you've built}]
A mountain, named for the knots in it.
\end{grammarlens}

\subsection*{Your turn}
Say these aloud:

\begin{itemize}
\raggedright
  \item \emph{parvataḥ}
  \item it once more, slowly
  \item \emph{setuḥ}, then \emph{parvataḥ}
\end{itemize}

\begin{tcolorbox}[breakable,colback=teal!4,colframe=teal!35!black,title={Before you move on}]
What does \sk{पर्वतः} mean? (\textquotedblleft{}a mountain\textquotedblright{}.) Say it once more.
\end{tcolorbox}
\section[puram]{\sk{पुरम्} (puram) — a town, a walled town}

\label{lesson:SA-C33-town}



\begin{quote}
Before the new word: what did \emph{parvataḥ} mean?
\end{quote}

\subsection*{You'll want to know: \sk{पुरम्}}
\textbf{\sk{पुरम्}} (\emph{puram}) — \textquotedblleft{}a town, a walled town\textquotedblright{}.

You already have \textbf{\sk{नगरम्}} for a city. \textbf{\sk{पुरम्}} is the older word, and it meant a stronghold first: a place with a wall around it.

Its Greek cousin is \emph{polis}, behind \emph{police}, \emph{policy}, \emph{politics} and every \emph{-polis} on a map. Sanskrit's own \textbf{-\sk{पुरम्}} is on the map just as thickly, worn down to \emph{-pur}: \emph{Jaipur}, \emph{Nagpur}, \emph{Kanpur}, and \emph{Singapore}, which began as \emph{Siṁhapura}, the lion town. Neuter \textbf{-\sk{अम्}}.

\begin{grammarlens}[title={What you've built}]
A town, and the reason \emph{police} and \emph{Nagpur} are relatives.
\end{grammarlens}

\subsection*{Your turn}
Say these aloud:

\begin{itemize}
\raggedright
  \item \emph{puram}
  \item it once more, slowly
  \item \emph{parvataḥ}, then \emph{puram}
\end{itemize}

\begin{tcolorbox}[breakable,colback=teal!4,colframe=teal!35!black,title={Before you move on}]
What does \sk{पुरम्} mean? (\textquotedblleft{}a town, a walled town\textquotedblright{}.) Say it once more.
\end{tcolorbox}
\section[samudraḥ]{\sk{समुद्रः} (samudraḥ) — the sea}

\label{lesson:SA-C33-sea}



\begin{quote}
Before the last word of the chapter: what did \emph{puram} mean, and what did \emph{nadī} mean?
\end{quote}

\subsection*{You'll want to know: \sk{समुद्रः}}
\textbf{\sk{समुद्रः}} (\emph{samudraḥ}) — \textquotedblleft{}the sea\textquotedblright{}.

\textbf{\sk{समुद्रः}} is \emph{sam}, \textquotedblleft{}together\textquotedblright{}, joined to \emph{udra}, \textquotedblleft{}water\textquotedblright{}: the sea is the gathering of the waters, which is where a \textbf{\sk{नदी}} ends up.

That \emph{sam-} at the front is the one inside \textbf{\sk{संस्कृतम्}}, \textquotedblleft{}put together, made complete\textquotedblright{}, the name of the language you have been saying since you first said you speak it. And \emph{udra} has cousins everywhere: Greek \emph{hudōr} behind \emph{hydrant} and \emph{hydrogen}, English \emph{water}, Russian \emph{voda}, and the \emph{otter}, who is nothing more than the water animal.

\begin{grammarlens}[title={What you've built}]
Five pieces of country, and the place a river finishes.
\end{grammarlens}

\subsection*{Your turn}
Say these aloud:

\begin{itemize}
\raggedright
  \item \emph{samudraḥ}
  \item it once more, slowly
  \item all five in order, then say where a \emph{nadī} ends
\end{itemize}

\begin{tcolorbox}[breakable,colback=teal!4,colframe=teal!35!black,title={Before you move on}]
What does \sk{समुद्रः} mean? (\textquotedblleft{}the sea\textquotedblright{}.) Say it once more.
\end{tcolorbox}
